
\begin{proposition}[The analytic support comparison used here]\label{cdp:st:C.2.MH.1}
Let $U$ be a smooth complex threefold, let $S\subset U$ be a
compact divisor, and suppose that an embedded resolution
$r:U'\to U$ is a sequence of blowups with smooth compact
centres in $S$, with $D=(r^{-1}S)_{\rm red}$ a simple
normal crossings divisor whose components are projective.
Let $V$ be a finite-monodromy rational local system on $U$.
Suppose that the cohomology of the punctured sufficiently
small fibre tubes around $S$, with coefficients in $V$,
vanishes in every degree. Then the mixed Hodge structure
on $H^k(S,V)$ obtained from a projective hyperresolution
is pure of weight $k$.

In this assertion the support-to-restriction morphism is
the actual morphism
\[
 R\Gamma_D(U',r^*V)\longrightarrow
 \underset{D\subset W}{\mathop{\rm colim}}\,
 R\Gamma(W,r^*V)=R\Gamma(D,r^*V).
 \tag{C.MH.1}\label{cdp:eq:C.MH.1}
\]
Neither a K\"ahler metric on $U$ nor a projective
compactification of $U$ is needed for this comparison.
\end{proposition}
\begin{proof}
First take constant rational coefficients. Write $D_I$
for the intersection of the components indexed by a
nonempty set $I$, and $D^{[p]}=\coprod_{|I|=p+1}D_I$.
All these spaces are smooth and projective. The ordinary
\v Cech de Rham complex on $D^{[\bullet]}$, with its
\v Cech weight filtration and its differential-form
filtration, is therefore a cohomological mixed Hodge
complex $K(D)$ for $\mathbb Q_D$. Its cohomology in
degree $k$ has weights at most $k$.

Here is the precise analytic construction of the other
complex and of the arrow. For each stratum $C=D_I$, put
\[
 K(C\log D)_{\mathbb C}
   =\Omega_{U'}^\bullet(\log D)\otimes\OO_C.
 \tag{C.MH.2}\label{cdp:eq:C.MH.2}
\]
The tensor product in \eqref{cdp:eq:C.MH.2} is a complex: in local
normal crossings coordinates, the ideal of $C$ is
preserved by the logarithmic differential. Endow it with
the quotient residue filtration and the filtration by
form degree. The quotient map from the restriction of
$\Omega_{U'}^\bullet(\log D)$ is a filtered
quasi-isomorphism to \eqref{cdp:eq:C.MH.2}. Its rational realization is
$i_C^*Rj_*\mathbb Q_{U'\setminus D}$, where $j$ is the
open inclusion. To see the comparison, apply the
iterated Poincar\'e residue: its residue-degree-$a$
quotient is a direct sum of the de Rham complexes of
$D_J\cap C$, shifted by $[-a]$ and with Tate twist
$(-a)$. The ideal and its differentials restrict these
forms to precisely $D_J\cap C$. This is the local
coordinate proof of \cite[Section~3, Theorem~1 and
Corollary~2]{Steenbrink1999}; it uses holomorphic
coordinates and the projectivity, or compact K\"ahler
property, of the displayed strata. It does not use a
compactification of their ambient neighbourhood.

Totalizing \eqref{cdp:eq:C.MH.2} over $D^{[\bullet]}$ gives the
cohomological mixed Hodge complex $K(D\log D)$, with
rational realization $i_D^*Rj_*\mathbb Q_{U'\setminus D}$.
The natural inclusion $K(D)\to K(D\log D)$ is a
morphism of mixed Hodge complexes. Its mixed cone,
shifted by $[-1]$, computes cohomology with support in
$D$; the cone projection is exactly \eqref{cdp:eq:C.MH.1}. This is
the excision comparison of
\cite[Section~4, equation~(3)]{Steenbrink1999}, applied
on the compact divisor. All complexes just used are
defined on that divisor, and all purity checks are on
its projective strata.

For completeness, the weight bound on the source can
be checked without invoking an invariant-cycle theorem.
The closed-support Mayer--Vietoris complex, together
with the holomorphic Thom isomorphisms, gives
\[
 E_1^{-p,q}=
 \bigoplus_{|I|=p+1}
 H^{q-2p-2}(D_I,\mathbb Q)(-p-1)
 \ \Longrightarrow\ H_D^{q-p}(U',\mathbb Q),
 \qquad p\ge0.
 \tag{C.MH.3}\label{cdp:eq:C.MH.3}
\]
The first differential is the alternating sum of Gysin
maps for $D_I\hookrightarrow D_{I\setminus\{i\}}$.
The comparison with the preceding mixed cone can be
verified by induction on the number of components.
For one smooth divisor the residue exact sequence
identifies
\[
 \operatorname{Cone}\bigl[\Omega_{U'}^\bullet
       \longrightarrow\Omega_{U'}^\bullet(\log D_i)\bigr][-1]
       \simeq\Omega_{D_i}^\bullet[-2](-1),
\]
where $(-1)$ records both the shift of the Hodge
filtration by one and the rational $2\pi i$ residue.
For a union $A\cup B$, use the actual support triangle
$R\Gamma_{A\cap B}\to R\Gamma_A\oplus R\Gamma_B
\to R\Gamma_{A\cup B}$. Iterated residues commute
with its inclusion maps, with the alternating signs
fixed by the ordering of the omitted coordinates.
Induction identifies its totalization with the same
mixed cone. These comparisons agree on overlaps,
since they are the residue maps for the holomorphic
divisor line bundles. Thus \eqref{cdp:eq:C.MH.3} compares the actual support
complex, including its arrow to $K(D)$, rather than
only its unfiltered vector spaces. Its displayed term
has pure weight $q$. In total degree $k=q-p$ this is
weight $k+p\ge k$. It follows that $H_D^k(U',\mathbb Q)$
has weights at least $k$.

Finite coefficients are handled by taking a finite
\'{e}tale cover on which $V$ is constant and then a
rational direct summand of a finite direct sum of its
regular representation.
The inverse images of the $D_I$ are still projective:
a finite holomorphic \'{e}tale cover of a
projective manifold is projective, since the pullback
of an ample bundle is positive. The deck action acts
on the displayed mixed Hodge complexes. Rational
idempotents consequently give exactly the same
comparison and weight bounds with coefficients in $V$.

By the assumed punctured-tube acyclicity, the support
long exact sequence makes \eqref{cdp:eq:C.MH.1} an isomorphism on
cohomology. Properness of $r$ identifies the punctured
tubes upstairs and downstairs, since its centres lie
in $S$. The map is a morphism of mixed Hodge structures,
whose source has weights at least $k$ and whose target
has weights at most $k$. Strictness therefore proves
that $H^k(D,r^*V)$ is pure of weight $k$.

We finally make the comparison with the original
support explicit. The proper degree-one map between
the smooth ambient manifolds has the trace map, defined
by oriented Poincar\'e duality and proper pushforward,
and the projection formula gives
\[
 r_*r^*=\operatorname{id}:H^k(U,V)\longrightarrow H^k(U,V).
 \tag{C.MH.4}\label{cdp:eq:C.MH.4}
\]
The same identity holds on every smaller ambient
neighbourhood. Passing to the direct limit of these
identities gives an injection
\[
                  r^*:H^k(S,V)\hookrightarrow H^k(D,r^*V).
 \tag{C.MH.5}\label{cdp:eq:C.MH.5}
\]
The identifications in this limit are the usual
continuity identifications for a compact triangulable
subset. They can equally be made using cofinal
triangulated regular neighbourhoods retracting onto
$S$ and $D$. Properness ensures that inverse images of
neighbourhoods are cofinal; thus the map in \eqref{cdp:eq:C.MH.5}
is the pullback by $D\to S$.

This last map is a mixed Hodge morphism as well. One
may construct the projective hyperresolution of $S$
from its projective normalization resolution and
resolutions of the conductor and exceptional curves.
To compare it with $D^{[\bullet]}$, resolve their
fibre products over $S$. Each such fibre product is a
closed analytic subspace of a product of projective
spaces already occurring in the two resolutions;
Chow's theorem makes it algebraic and projective,
and its resolution is projective. The resulting
common hyperresolution realizes \eqref{cdp:eq:C.MH.5} by pullbacks
between projective strata. Cohomological descent then
identifies this with the stated mixed Hodge structures.
Injectivity into the pure target proves the assertion.
\end{proof}


